\chapter{Conclusão}
\label{cap:conclusão}
\section{Contribuições do trabalho}

No contexto da análise dos modelos de predição de óbitos no cenário da Covid-19 a medida que a vacinação ocorre, este trabalho mostrou através do mapeamento sistemático os principais modelos de redes neurais escolhidos para esse contexto, entre eles estão: 
\begin{itemize}
    \item SIRUCQHE (Susceptible - Infected - Removed - Unsusceptible - Contagious - Quarantined - Hospitalized - Extict) \cite{Scarabaggio};
    \item SEIAQFR (Susceptible - Exposed - Infected - Asymptomatic - Quarantined - Fatal - Recovered) \cite{Rakshit};
    \item SIR (Susceptible - Infected - Recovered) \cite{Chinchilla};
    \item DNN (Fully Connected Deep Neural Network) \cite{Alkhammash};
    \item LSTM (Long Short-Term Memory) \cite{Nikita} e \cite{Alkhammash};
    \item TM (Transformer Model) \cite{Alkhammash};
    \item RNN (Recurrent Neural Network) \cite{Nikita};
    \item CNN (Convolutional Neural Network) \cite{Nikita};
    \item ANN (Multi-Layer Perceptron Artificial Neural Network) \cite{DeCarvalho}.
\end{itemize} 

% SIRUCQHE (Suscetível-Infectado-Removido-Insusceptível-Contagioso-Em Quarentena-Hospitalizado-Extinto.) \cite{Scarabaggio}, SEIAQFR (Suscetível-Exposto-Infectado-Assintomático-Em Quarentena-Fatal-Recuperado) \cite{DeCarvalho}, SIR (Suscetível-Infectado-Recuperado) \cite{Chinchilla}, DNN (Rede Neural Profunda Totalmente Conectada) \cite{Alkhammash}, LSTM (Memória de Curto Longo Prazo) \cite{Nikita} e \cite{Alkhammash}, TM (Modelo Transformador) \cite{Alkhammash}, RNN (Rede Neural Recorrente) \cite{Nikita}, CNN (Rede Neural Convolucional) \cite{Nikita} e ANN (Rede Neural Artificial de Perceptron Multicamadas) \cite{Rakshit}.

Além disso, foi realizado um estudo de caso na região Centro-Oeste brasileiro, evidenciando através de gráficos as curvas de vacinação mensal a partir de janeiro de 2021 até junho de 2022 e os tipos de doses.

Por fim, foi implementado um algoritmo de predição baseado no modelo LSTM, o conjunto de dados que alimenta esse modelo foi retirado do \textit{OpenDataSUS} referente ao Estado de Goiás, e nele contém a quantidade de óbitos, vacinados e casos novos, todos mensalmente. O algoritmo apresentou um overfitting com causa não identificada.

\section{Trabalhos Futuros}

Para um próximo passo melhorar a qualidade do algortimo de predição utilizando o modelo LSTM e acrescentar mais algoritmos para análise estatística de óbitos.

\section{Disponibilidade dos Códigos e Dados}
Todos os algoritmos criados para este trabalho estão disponíveis no Github\cite{Github}, para o modelo LSTM próprio foi inspirado no algoritmo retirado do site \textit{Dynamic Optimization}\cite{LSTM}. E os dados que foram retirados do \textit{OpenDataSUS} são os da Campanha Nacional de Vacinação contra Covid-19\cite{OpenDataSus}, e do Registro de Ocupação Hospitalar COVID-19\cite{Open}. Além desses conjuntos de dados também foi retirado uma tabela já pré-processada do site Coronavírus Brasil\cite{Coronavirus}.
